% Copyright (c) 2026 Selina Cheggour. Tous droits réservés. Voir le fichier LICENSE.
% Annexe : les outils expliqués (babel, py3langid, expressions régulières et nettoyage, centrage par langue).

\bandeaupartie{Annexe : les outils expliqués}{babel, py3langid, les expressions régulières du nettoyage,\\et le centrage par langue expliqué avec des notes de classe.}

\section{babel : le nom des langues en clair}

\begin{stepbox}[title={\faLanguage\quad Ce que c'est}]
\texttt{babel} est une bibliothèque Python d'\textbf{internationalisation} : pour chaque langue d'affichage,
elle connaît les noms des langues et des pays, les formats de dates, de nombres et de monnaies. Elle les
tire du \textbf{CLDR} (\emph{Common Locale Data Repository}), la base de référence du consortium Unicode,
celle qu'utilisent aussi les navigateurs et les systèmes d'exploitation.
\end{stepbox}

\begin{stepbox}[title={\faCogs\quad Ce qu'elle fait dans le notebook}]
Une seule chose : écrire en clair le nom des langues détectées. Le détecteur de langue (ci-dessous) ne
renvoie qu'un \textbf{code}, comme \texttt{fr}, \texttt{nl} ou \texttt{ary}. babel le traduit en nom
français :
\begin{lstlisting}[escapeinside={(*}{*)}]
from babel import Locale
NOMS_CLDR = Locale("fr").languages      # un dictionnaire : code -> nom en français
NOMS_CLDR["nl"]                          # 'néerlandais'
NOMS_CLDR["ary"]                         # 'arabe marocain'
\end{lstlisting}
Quelques codes rares que le CLDR ne nomme pas sont complétés à la main dans \texttt{COMPLEMENT\_LANGUES}
(par exemple \texttt{zxx}, \og sans contenu linguistique \fg). Le nom affiché est donc, dans l'ordre : le
complément s'il existe, sinon le nom donné par babel, sinon le code lui-même.
\end{stepbox}

\begin{tcbraster}[raster columns=2, raster equal height, raster column skip=8pt, raster every box/.style={unbreakable}]
\begin{conceptbox}[title={\faDownload\ Installée en 0.2}, unbreakable]
\small Elle n'est pas dans l'image Kaggle par défaut : la cellule 0.2 l'installe avec \texttt{py3langid}.
Si elle manque, rien ne plante : les langues s'affichent par leur code.
\end{conceptbox}
\begin{conceptbox}[title={\faLock\ Ne touche pas aux plaintes}, unbreakable]
\small Elle ne lit aucun texte et n'envoie rien sur internet. Elle ne joue aucun rôle dans les vecteurs, les
groupes ou les alertes : elle sert seulement à l'affichage.
\end{conceptbox}
\end{tcbraster}

%%%%%%%%%%%%%%%%%%%%%%%%%%%%%%%%%%%%%%%%%%%%%%%%%%%%%%%%%%%%%%%%%%%%%%%%%%%%%%%%%%%%%%%%%%%%%%%%%%%%%%%%
\section{py3langid : quelle langue pour chaque plainte}

\begin{stepbox}[title={\faSearch\quad Ce que c'est}]
\texttt{py3langid} est un \textbf{détecteur de langue} : on lui donne un texte, il renvoie le code de la
langue la plus probable. C'est la version modernisée de \texttt{langid.py}. Son modèle, un petit modèle
probabiliste (Bayes naïf), a appris sur des textes de plus d'une centaine de langues quelles
\textbf{suites de caractères} sont typiques de chacune : \og ijk \fg{} en néerlandais, \og sch \fg{} en
allemand, \og eau \fg{} en français. Il est livré avec le paquet et tourne sur le processeur, sans internet
et sans carte graphique, en une fraction de seconde pour toute la base.
\end{stepbox}

\begin{stepbox}[title={\faCogs\quad Ce qu'il fait dans le notebook (exercice 1.4)}]
\begin{lstlisting}[escapeinside={(*}{*)}]
import py3langid
py3langid.classify("Bonjour, le chocolat avait un go(*\^u*)t de savon.")
# ('fr', -209.3)  : le code de la langue, puis un score
plaintes["langue"] = [py3langid.classify(t)[0] if t else "?" for t in plaintes.texte]
\end{lstlisting}
\texttt{classify} renvoie deux valeurs ; on garde la première, le code (\texttt{[0]}). Le score est une
\textbf{log-probabilité} : plus il est grand, plus le modèle est sûr, mais ce n'est pas un pourcentage. Un
texte vide reçoit \texttt{"?"}.
\doisvoir{une langue pour chaque plainte, et une répartition avec quelques langues principales.}
\end{stepbox}

\begin{piege}
\textbf{La langue n'est pas le pays.} L'anglais, le français et l'allemand couvrent chacun plusieurs pays.
\textbf{Et le détecteur se trompe sur les textes très courts} ou pleins de chiffres et de références : une
langue qui ne compte qu'une ou deux plaintes est souvent une erreur de détection, pas un vrai pays de plus.
\end{piege}

%%%%%%%%%%%%%%%%%%%%%%%%%%%%%%%%%%%%%%%%%%%%%%%%%%%%%%%%%%%%%%%%%%%%%%%%%%%%%%%%%%%%%%%%%%%%%%%%%%%%%%%%
\section{Les expressions régulières du nettoyage (exercice 1.2)}

\subsection{re, c'est quoi ?}

\texttt{re} est le module de Python pour les \textbf{expressions régulières} (\emph{regular expressions}).
Une expression régulière est un \textbf{motif} qui décrit une \emph{forme} de texte, au lieu d'un texte
exact. Au lieu de chercher le lien précis \texttt{https://exemple.io}, on décrit la forme de tous les liens :
\og commence par http, peut-être un s, puis ://, puis des caractères sans espace \fg.

\subsection{Comment se lit re.sub}

\texttt{re.sub} veut dire \emph{substitute}, remplacer. Elle prend toujours trois choses, dans cet ordre :
\begin{center}
\begin{tikzpicture}[b/.style={draw=#1, fill=#1!10, rounded corners=4pt, font=\small, align=center, minimum height=1cm, text width=3.6cm}]
\node[b=sectioncolor] (m) {\textbf{1. le motif}\\ce qu'on cherche};
\node[b=accentblue, right=0.4cm of m] (r) {\textbf{2. le remplacement}\\par quoi on le remplace};
\node[b=darkblue, right=0.4cm of r] (t) {\textbf{3. le texte}\\où on cherche};
\node[above=0.15cm of r, font=\ttfamily\small] {re.sub( motif , remplacement , texte )};
\end{tikzpicture}
\end{center}
Elle renvoie le texte modifié. Exemple : \verb|re.sub(r"\*{2,}", " ", texte)| se lit \og dans
\texttt{texte}, remplace chaque suite de deux étoiles ou plus par une espace \fg.

\begin{tcbraster}[raster columns=2, raster equal height, raster column skip=8pt, raster every box/.style={unbreakable}]
\begin{conceptbox}[title={Le \texttt{r} devant le motif}, unbreakable]
\small \verb|r"..."| est un texte \og brut \fg{} : Python garde les barres \verb|\| telles quelles et les
passe au motif.
\end{conceptbox}
\begin{conceptbox}[title={\texttt{texte}, à la fin}, unbreakable]
\small C'est la variable qui contient la plainte à nettoyer. Le résultat de chaque ligne remplace
\texttt{texte}, et la ligne suivante repart de lui.
\end{conceptbox}
\end{tcbraster}

\subsection{Les briques pour lire les motifs}

\begin{center}
\renewcommand{\arraystretch}{1.3}
\rowcolors{2}{notebg}{white}
\begin{tabular}{@{}>{\centering\arraybackslash}p{2.2cm}p{11.8cm}@{}}
\rowcolor{titlebg}\color{white}\textbf{symbole} & \color{white}\textbf{veut dire} \\
\verb|^| & le tout début du texte \\
\verb|\s| & un \og blanc \fg{} : espace, tabulation, retour à la ligne \\
\verb|\S| & tout sauf un blanc \\
\verb|*| & 0 fois ou plus le symbole d'avant \\
\verb|+| & 1 fois ou plus le symbole d'avant \\
\verb|?| & 0 ou 1 fois le symbole d'avant : il est facultatif \\
\verb|{2,}| & 2 fois ou plus le symbole d'avant \\
\verb|\*| & une vraie étoile (la barre empêche de la lire comme \og 0 fois ou plus \fg) \\
\verb|[.!?]| & un seul caractère parmi ceux de la liste \\
\verb|[^\W\d_]| & une lettre, de n'importe quel alphabet (ni chiffre, ni ponctuation, ni \_) \\
\verb|(...)| et \verb|\1| & on capture un morceau, puis on le remet tel quel dans le remplacement \\
\verb|(?=...)| & regarde ce qui suit, sans le consommer ni le modifier \\
\end{tabular}
\end{center}

\subsection{Les cinq étapes de nettoyer, une par une}

\begin{lstlisting}[escapeinside={(*}{*)}]
def nettoyer(texte):
    texte = re.sub(r"^\s*EXTERNAL", "", texte)               # 1
    texte = re.sub(r"https?://\S+", " ", texte)               # 2
    texte = re.sub(r"\*{2,}", " ", texte)                     # 3
    texte = re.sub(r"([.!?;:,(*\textcolor{accentblue}{\textsf{\scriptsize virgule arabe}}*)])(?=[^\W\d_])", r"\1 ", texte) # 4
    return re.sub(r"\s+", " ", texte).strip()                 # 5
\end{lstlisting}

\begin{center}
\renewcommand{\arraystretch}{1.35}
\rowcolors{2}{notebg}{white}
\begin{tabular}{@{}>{\bfseries\color{sectioncolor}}c p{6.4cm}p{6.6cm}@{}}
\rowcolor{titlebg}\color{white}\# & \color{white}\textbf{ce que fait le motif} & \color{white}\textbf{exemple} \\
1 & au début du texte (\verb|^|), des blancs ou rien, puis le mot EXTERNAL : supprimé & \texttt{EXTERNALBonjour} $\to$ \texttt{Bonjour} \\
2 & \texttt{http}, un \texttt{s} facultatif, \texttt{://}, puis tout jusqu'au premier blanc : le lien entier devient une espace & \texttt{voir https://exemple.io/aB3} $\to$ \texttt{voir \textvisiblespace} \\
3 & deux étoiles ou plus (les produits masqués) : une espace & \texttt{acheté ******* hier} $\to$ \texttt{acheté \textvisiblespace\textvisiblespace\textvisiblespace hier} \\
4 & une ponctuation suivie directement d'une lettre : on garde la ponctuation (\verb|\1|) et on ajoute une espace & \texttt{Bonjour,J'ai} $\to$ \texttt{Bonjour, J'ai} ; \texttt{3.5} ne bouge pas \\
5 & les espaces simples : chaque groupe de blancs devient une seule espace, puis \texttt{.strip()} enlève les blancs du début et de la fin & \texttt{acheté \textvisiblespace\textvisiblespace\textvisiblespace hier} $\to$ \texttt{acheté hier} \\
\end{tabular}
\end{center}
{\small\color{slate}\textit{Le signe \textvisiblespace{} montre une espace.}}

\subsection{L'étape 5 et strip, en détail}

\begin{stepbox}[title={5\quad Des espaces simples}]
La dernière ligne fait deux choses, de gauche à droite.
\begin{enumerate}[nosep]
\item \verb|re.sub(r"\s+", " ", texte)| : \verb|\s+| veut dire \og un blanc ou plus, d'affilée \fg. Chaque
  groupe de blancs, qu'il compte 1 ou 10 espaces, tabulations ou retours à la ligne, devient \textbf{une
  seule espace}.
\item \texttt{.strip()} n'est pas du \texttt{re} : c'est une méthode de tout texte Python. Elle enlève les
  blancs au \textbf{début} et à la \textbf{fin}, jamais au milieu :
  \verb|"  Bonjour le monde  ".strip()| donne \verb|"Bonjour le monde"|.
\end{enumerate}
Cette étape est nécessaire parce que les étapes d'avant \textbf{créent} des espaces en trop :
\begin{lstlisting}[escapeinside={(*}{*)}]
départ      : "EXTERNALBonjour, j'ai acheté ******* hier"
après 1     : "Bonjour, j'ai acheté ******* hier"
après 3     : "Bonjour, j'ai acheté   hier"     (trois espaces)
après 5     : "Bonjour, j'ai acheté hier"       (une seule)
\end{lstlisting}
Un lien en fin de message devient lui aussi une espace finale, que \texttt{.strip()} retire.
\image{après avoir découpé des étiquettes dans une lettre, on referme les trous pour que le texte se lise d'une traite.}
\end{stepbox}

%%%%%%%%%%%%%%%%%%%%%%%%%%%%%%%%%%%%%%%%%%%%%%%%%%%%%%%%%%%%%%%%%%%%%%%%%%%%%%%%%%%%%%%%%%%%%%%%%%%%%%%%
\section{Le centrage par langue, avec des notes de classe (exercice 2.3)}

\subsection{Deux professeurs, deux façons de noter}

\begin{tcbraster}[raster columns=2, raster equal height, raster column skip=8pt, raster every box/.style={unbreakable}]
\begin{bloc}[sectioncolor]{\faUserTie\quad La classe du professeur dur}
\begin{center}
\renewcommand{\arraystretch}{1.2}
\rowcolors{2}{notebg}{white}
\begin{tabular}{lr}
\textbf{élève} & \textbf{note} \\
Ali & 6 \\
Bea & 8 \\
Chloé & 9 \\
Dan & 10 \\
\textbf{Léa} & \textbf{12} \\
\end{tabular}
\end{center}
Moyenne de la classe : $(6 + 8 + 9 + 10 + 12) / 5 = 45 / 5 = 9$.
\end{bloc}
\begin{bloc}[accentblue]{\faUserGraduate\quad La classe du professeur large}
\begin{center}
\renewcommand{\arraystretch}{1.2}
\rowcolors{2}{notebg}{white}
\begin{tabular}{lr}
\textbf{élève} & \textbf{note} \\
\textbf{Tom} & \textbf{12} \\
Uma & 14 \\
Vic & 15 \\
Wael & 16 \\
Zoé & 18 \\
\end{tabular}
\end{center}
Moyenne de la classe : $(12 + 14 + 15 + 16 + 18) / 5 = 75 / 5 = 15$.
\end{bloc}
\end{tcbraster}

\subsection{Même note, deux histoires opposées}

\begin{tcbraster}[raster columns=2, raster equal height, raster column skip=8pt, raster every box/.style={unbreakable}]
\begin{kpibox}[title={Léa : $12 - 9 = +3$}, unbreakable]
\small Léa a 3 points \textbf{de plus} que ce que son professeur donne d'habitude. C'est la
\textbf{meilleure} de sa classe.
\end{kpibox}
\begin{kpibox}[title={Tom : $12 - 15 = -3$}, unbreakable]
\small Tom a 3 points \textbf{de moins} que ce que son professeur donne d'habitude. C'est le
\textbf{dernier} de sa classe.
\end{kpibox}
\end{tcbraster}

Les deux ont 12. Le $+3$ et le $-3$ montrent la différence ; le 12 brut la cachait. Le
\textbf{centrage}, c'est juste ça : \textbf{remplacer la note par l'écart à la moyenne de son propre
professeur}. On ne compare plus des notes, on compare ce que chaque élève a de plus ou de moins que sa
classe.

\subsection{Le lien avec les plaintes}

\begin{center}
\renewcommand{\arraystretch}{1.35}
\rowcolors{2}{notebg}{white}
\begin{tabular}{@{}p{4.6cm}p{9.4cm}@{}}
\rowcolor{titlebg}\color{white}\textbf{dans la classe} & \color{white}\textbf{dans le notebook} \\
le professeur & la \textbf{langue} de la plainte \\
la note d'un élève & le \textbf{plongement} $\mathbf{e}_i$ de la plainte : 1\,024 nombres au lieu d'un seul \\
la moyenne de la classe & la moyenne $\boldsymbol{\mu}_\ell$ des plongements de toutes les plaintes de la langue $\ell$ \\
l'écart $+3$ ou $-3$ & $\mathbf{e}_i - \boldsymbol{\mu}_{\ell(i)}$ : ce que la plainte dit \textbf{de plus} que les autres plaintes de sa langue \\
\end{tabular}
\end{center}

Un plongement mélange deux choses : ce que dit la plainte (le sujet) et la façon dont sa langue
s'exprime (la \og façon de noter \fg{} du professeur). Deux plaintes allemandes sur des sujets différents
se ressemblent un peu, simplement parce qu'elles sont en allemand. En retirant à chaque plainte la
moyenne de sa langue, on retire cette part commune, et il reste le sujet : une plainte française et une
plainte allemande sur le goût de savon se retrouvent alors côte à côte.

\begin{tcbraster}[raster columns=2, raster equal height, raster column skip=8pt, raster every box/.style={unbreakable}]
\begin{conceptbox}[title={\faRulerCombined\ Puis on renorme}, unbreakable]
\small Après la soustraction, on divise par la longueur du vecteur (le dénominateur de la formule de
l'exercice 2.3), pour que toutes les plaintes aient la même longueur, 1.
\end{conceptbox}
\begin{conceptbox}[title={\faUsers\ Les petites classes}, unbreakable]
\small Une moyenne sur deux élèves ne dit rien du professeur. De même, une langue de moins de 10 plaintes
reçoit la moyenne de \textbf{toutes} les plaintes, pas la sienne.
\end{conceptbox}
\end{tcbraster}

% Annexe : cosinus, similarité et distance (exercices 2.1 et 3.2), en images.

\section{Cosinus : similarité et distance (exercices 2.1 et 3.2)}

\subsection{Le cosinus regarde l'angle entre deux flèches}

Chaque plainte, et chaque centre de groupe, est une \textbf{flèche} (un vecteur). Le \textbf{cosinus} ne
regarde pas la longueur des flèches, seulement l'\textbf{angle} entre elles : il dit si elles pointent dans
la même direction, donc si les deux textes parlent de la même chose.

\begin{center}
\begin{tikzpicture}[fl/.style={-{Stealth[length=3mm]}, line width=1.6pt}]
\foreach \dx/\ang/\cosv/\dv/\titre/\coul in {0/8/1/0/même direction/accentblue, 5.4/60/0{,}5/0{,}5/à moitié/ambre, 10.8/90/0/1/rien en commun/slate} {
  \begin{scope}[xshift=\dx cm]
  \fill[palblue!40, rounded corners=6pt] (-0.6,-1.7) rectangle (4.2,2.9);
  \node[font=\small\bfseries, color=\coul] at (1.8,2.55) {\titre};
  \draw[fl, color=sectioncolor] (0,0) -- (3.2,0);
  \draw[fl, color=\coul] (0,0) -- (\ang:2.3);
  \ifnum\ang>10 \draw[color=\coul, thick] (0.75,0) arc (0:\ang:0.75); \node[font=\scriptsize, color=\coul] at ({\ang/2}:1.1) {\ang\textdegree}; \fi
  \node[font=\small, anchor=west] at (-0.4,-0.65) {$\cos = \cosv$};
  \node[font=\small, anchor=west] at (-0.4,-1.25) {distance $1 - \cos = \dv$};
  \end{scope}
}
\end{tikzpicture}
\end{center}
{\small\color{slate}\textit{Flèche rouge : une plainte ; flèche colorée : une autre plainte. Seul l'angle compte.}}

\subsection{Deux façons de dire la même chose}

\begin{tcbraster}[raster columns=2, raster equal height, raster column skip=8pt, raster every box/.style={unbreakable}]
\begin{conceptbox}[title={\faHeart\ La similarité cosinus}, unbreakable]
\small \textbf{Grande} quand les textes se ressemblent : 1 pour la même direction, 0 pour rien en commun.
C'est elle qu'on lit dans les seuils du TP (\og cosinus $> 0{,}4$ \fg, voisine \og au-delà de 0,9 \fg).
\end{conceptbox}
\begin{conceptbox}[title={\faRulerHorizontal\ La distance cosinus}, unbreakable]
\small \textbf{Petite} quand les textes se ressemblent : $d = 1 - \cos$. C'est ce que demandent certaines
fonctions de scikit-learn, comme la fusion des jumeaux (\texttt{metric="cosine"}).
\end{conceptbox}
\end{tcbraster}

\begin{center}
\begin{tikzpicture}[x=1.35cm]
% règle du cosinus (en haut) et de la distance (en bas), alignées
\shade[left color=accentblue!70, right color=palblue] (0,0) rectangle (10,0.45);
\draw[slate] (0,0) rectangle (10,0.45);
\foreach \x/\c/\d in {0/1/0, 2/0{,}8/0{,}2, 4/0{,}6/0{,}4, 6/0{,}4/0{,}6, 8/0{,}2/0{,}8, 10/0/1} {
  \draw[slate] (\x,0.45) -- (\x,0.6); \node[above, font=\small] at (\x,0.6) {\c};
  \draw[slate] (\x,0) -- (\x,-0.15); \node[below, font=\small] at (\x,-0.15) {\d};
}
\node[left, font=\small\bfseries, color=accentblue] at (-0.25,0.85) {cosinus};
\node[left, font=\small\bfseries, color=darkblue] at (-0.25,-0.45) {distance};
% seuil de la fusion des jumeaux
\draw[very thick, dashed, color=sectioncolor] (6,-0.05) -- (6,0.5);
\draw[very thick, dashed, color=sectioncolor] (6,1.12) -- (6,1.38);
\draw[very thick, dashed, color=sectioncolor] (6,-0.62) -- (6,-0.95);
\node[font=\footnotesize\bfseries, color=sectioncolor, align=center] at (6,1.75) {seuil par défaut\\$\cos = 0{,}4$, $d = 0{,}6$};
\draw[-{Stealth[length=3mm]}, very thick, color=accentblue] (5.85,-1.15) -- (0.2,-1.15)
  node[midway, below, font=\footnotesize] {on fusionne : $\cos > 0{,}4$, donc $d < 0{,}6$};
\draw[-{Stealth[length=3mm]}, very thick, color=slate] (6.15,-1.15) -- (9.8,-1.15)
  node[midway, below, font=\footnotesize] {on ne fusionne pas};
\end{tikzpicture}
\end{center}

\begin{lire}
\textbf{Comment lire la règle.} Le haut donne le cosinus, le bas la distance au même endroit : les deux
nombres d'une même graduation font toujours 1 à eux deux. \og Fusionner deux groupes si le cosinus de leurs
centres dépasse 0,4 \fg{} et \og fusionner si leur distance est inférieure à 0,6 \fg{} désignent donc
exactement la même zone, à gauche du trait. Dans le code, on règle le cosinus et la conversion se fait toute
seule : \texttt{distance\_threshold = 1 - cosinus\_min}. Le réglage automatique (exercice 3.3) peut monter le
seuil à $\cos > 0{,}7$, soit $d < 0{,}3$ : il ne fusionne alors que des groupes très proches.
\end{lire}

\image{deux personnes qui montrent une direction du doigt : peu importe la longueur du bras, on regarde
seulement si elles montrent le même endroit.}

% Copyright (c) 2026 Selina Cheggour. Tous droits réservés. Voir le fichier LICENSE.
% Annexe : les méthodes eom, leaf (HDBSCAN) et Ward (plan B), exercices 3.1 et 3.3.

\newpage
\section{Trois façons de faire les tas : eom, leaf et Ward (exercices 3.1 et 3.3)}

\subsection{D'abord, l'arbre de HDBSCAN}

HDBSCAN regarde les plaintes comme des îles sur une carte, et fait \textbf{baisser le niveau de la mer}
petit à petit. Au début, tout est une seule grande île. Quand l'eau monte entre deux zones moins denses,
l'île se coupe en deux, puis encore en deux\ldots{} On obtient un \textbf{arbre} : chaque branche est un
groupe possible, et sa \textbf{longueur} dit combien de temps il a tenu avant de se couper. Une branche de
moins de \texttt{taille\_min} plaintes ne compte pas : ses plaintes deviennent du bruit (étiquette $-1$).

Il reste à choisir \textbf{quelles branches garder}. C'est là qu'interviennent les deux méthodes.

\begin{center}
\begin{tikzpicture}[x=0.62cm, y=0.62cm,
  br/.style={line width=#1, color=slate!70},
  pris/.style={line width=#1, color=accentblue},
  lab/.style={font=\scriptsize, color=slate}]
\foreach \dx/\titre/\coul in {0/{\texttt{eom} : les branches qui durent}/accentblue, 12.5/{\texttt{leaf} : les bouts de l'arbre}/ambre} {
\begin{scope}[xshift=\dx cm]
  \node[font=\small\bfseries, color=\coul] at (5,8.4) {\titre};
  % tronc
  \draw[br=5pt] (5,7.6) -- (5,6.5);
  \draw[br=2pt] (2.5,6.5) -- (7.5,6.5);
  % branche gauche A : longue, puis coupée en deux petits bouts
  \draw[br=4pt] (2.5,6.5) -- (2.5,2.6);
  \draw[br=1.5pt] (1.5,2.6) -- (3.5,2.6);
  \draw[br=2pt] (1.5,2.6) -- (1.5,1.4);
  \draw[br=2pt] (3.5,2.6) -- (3.5,1.6);
  % branche droite B : courte, puis coupée en deux longues branches
  \draw[br=4pt] (7.5,6.5) -- (7.5,5.2);
  \draw[br=1.5pt] (6.3,5.2) -- (8.7,5.2);
  \draw[br=3pt] (6.3,5.2) -- (6.3,0.8);
  \draw[br=3pt] (8.7,5.2) -- (8.7,1.2);
  \node[lab] at (2.5,7.0) {A}; \node[lab] at (7.5,7.0) {B};
  \node[lab, below] at (1.5,1.4) {A1}; \node[lab, below] at (3.5,1.6) {A2};
  \node[lab, below] at (6.3,0.8) {B1}; \node[lab, below] at (8.7,1.2) {B2};
\end{scope}
}
% eom : A (longue) et B1, B2 (longues) ; leaf : les quatre bouts
\draw[pris=4pt, color=accentblue] (2.5,6.5) -- (2.5,2.6);
\draw[pris=3pt, color=accentblue] (6.3,5.2) -- (6.3,0.8);
\draw[pris=3pt, color=accentblue] (8.7,5.2) -- (8.7,1.2);
\begin{scope}[xshift=12.5cm]
\draw[pris=2pt, color=ambre] (1.5,2.6) -- (1.5,1.4);
\draw[pris=2pt, color=ambre] (3.5,2.6) -- (3.5,1.6);
\draw[pris=3pt, color=ambre] (6.3,5.2) -- (6.3,0.8);
\draw[pris=3pt, color=ambre] (8.7,5.2) -- (8.7,1.2);
\end{scope}
\node[font=\footnotesize, text width=6.6cm, align=center] at (5,-0.6) {3 groupes : A entier, B1, B2};
\node[font=\footnotesize, text width=6.6cm, align=center] at (17.5,-0.6) {4 groupes : A1, A2, B1, B2};
\end{tikzpicture}
\end{center}
{\small\color{slate}\textit{Arbre inventé. En couleur, les branches gardées : chaque plainte ne peut être que dans une seule.}}

\begin{tcbraster}[raster columns=2, raster equal height, raster column skip=8pt, raster every box/.style={unbreakable}]
\begin{conceptbox}[title={\faTree\ \texttt{eom} : \og excess of mass \fg}, unbreakable]
\small Garde les branches qui \textbf{durent le plus longtemps}, en tenant compte de leur taille. Si une
branche dure plus que ses deux sous-branches réunies, on la garde entière (A) ; sinon on garde les
sous-branches (B1 et B2). Résultat : \textbf{moins de groupes, plus gros et plus solides}. C'est le choix
par défaut de HDBSCAN.
\end{conceptbox}
\begin{conceptbox}[title={\faLeaf\ \texttt{leaf} : les feuilles}, unbreakable]
\small Garde tous les \textbf{bouts de l'arbre}, les plus petites branches qui ne se coupent plus.
Résultat : \textbf{plus de groupes, plus petits et plus homogènes}. Utile quand on cherche des
problèmes fins, au risque de couper un même problème en morceaux.
\end{conceptbox}
\end{tcbraster}

\image{une famille qui partage un gâteau : \texttt{eom} coupe les grosses parts qui ont du sens (le côté
chocolat, le côté fraise), \texttt{leaf} coupe en toutes petites parts.}

\subsection{Le plan B : Ward}

Si aucun des 18 réglages de HDBSCAN n'est admis, c'est que les plaintes ne forment pas d'îles nettes.
On change alors de méthode : la classification de \textbf{Ward} part dans l'autre sens. Au début, chaque
plainte est seule dans son groupe ; à chaque étape, on \textbf{colle les deux groupes} dont la fusion
mélange le moins les plaintes, et on s'arrête quand il reste $k$ groupes.

\noindent\begin{minipage}{\linewidth}
\begin{center}
\begin{tikzpicture}[x=0.95cm, y=0.95cm, pt/.style={circle, fill=#1, inner sep=2.2pt},
  fl/.style={-{Stealth[length=2.5mm]}, thick, color=slate}]
\foreach \dx/\titre in {0/{départ : 6 groupes}, 4.6/{étape suivante : 4}, 9.2/{arrêt à $k = 2$}} {
  \node[font=\footnotesize\bfseries, color=darkblue] at (\dx+1.5,2.6) {\titre};
  \foreach \x/\y in {0.2/0.4, 0.8/0.9, 0.5/1.6, 2.2/0.5, 2.8/1.1, 2.4/1.8} \node[pt=sectioncolor] at (\dx+\x,\y) {};
}
% étape 2 : deux paires collées
\draw[accentblue, thick, rounded corners=6pt] (4.6+0.0,0.15) rectangle (4.6+1.05,1.15);
\draw[accentblue, thick, rounded corners=6pt] (4.6+1.95,0.25) rectangle (4.6+3.05,1.35);
% étape 3 : deux groupes
\draw[accentblue, very thick, rounded corners=8pt] (9.2-0.05,0.1) rectangle (9.2+0.85,1.9);
\draw[accentblue, very thick, rounded corners=8pt] (9.2+1.95,0.2) rectangle (9.2+3.05,2.1);
\draw[fl] (3.2,1.0) -- (4.3,1.0);
\draw[fl] (7.8,1.0) -- (8.9,1.0);
\end{tikzpicture}
\end{center}
{\small\color{slate}\textit{Exemple inventé avec 6 plaintes : à chaque étape, on colle les deux groupes les plus proches au sens de Ward.}}
\end{minipage}


\begin{formulabox}
\og Mélanger le moins \fg{} veut dire faire grandir le moins possible la \textbf{dispersion à l'intérieur des
groupes} :
\[
W = \sum_{h} \sum_{i \in h} \lVert \mathbf{z}_i - \bar{\mathbf{z}}_h \rVert^2 ,
\]
où $h$ parcourt les groupes, $\mathbf{z}_i$ est le vecteur de la plainte $i$ et $\bar{\mathbf{z}}_h$ la
moyenne du groupe $h$. Ward choisit à chaque étape la fusion qui augmente $W$ le moins.
\end{formulabox}

\begin{center}
\renewcommand{\arraystretch}{1.3}
\rowcolors{2}{notebg}{white}
\begin{tabular}{@{}>{\raggedright\arraybackslash}p{4.2cm}>{\raggedright\arraybackslash}p{5cm}>{\raggedright\arraybackslash}p{5cm}@{}}
\rowcolor{titlebg} & \color{white}\textbf{HDBSCAN (\texttt{eom} ou \texttt{leaf})} & \color{white}\textbf{Ward (plan B)} \\
nombre de groupes & trouvé tout seul & fixé d'avance : $k = 8, 12, 16, 20$ ou $30$ \\
plaintes isolées & laissées de côté ($-1$) & toujours rangées dans un groupe \\
ce qu'il cherche & des zones denses, de n'importe quelle forme & des groupes compacts, plutôt ronds \\
quand l'utiliser & des îles nettes dans les données & pas d'îles nettes : un continuum \\
\end{tabular}
\end{center}

\begin{piege}
\textbf{Ward range tout, même ce qui ne ressemble à rien.} Il n'y a plus de plaintes de côté : une
plainte bizarre finit quand même dans un groupe. C'est pourquoi on garde les mêmes règles d'admission
(pas de géant, assez de problèmes) et la fusion des jumeaux après Ward.
\end{piege}

\subsection{HDBSCAN et Ward : deux arbres, une seule différence qui compte}

Les deux méthodes sont \textbf{hiérarchiques}. Toutes les deux construisent un arbre de groupes emboîtés,
HDBSCAN en regardant la mer monter, Ward en collant les groupes deux à deux, puis décident où couper cet
arbre. Mais \textbf{seul HDBSCAN regarde la densité}. Ward ne la regarde jamais : il regarde la
\textbf{dispersion}. À chaque étape, il colle les deux groupes dont la fusion fait le moins grossir la
dispersion totale $W$, et il s'arrête quand il reste $k$ groupes.

\begin{tcbraster}[raster columns=2, raster equal height, raster column skip=8pt, raster every box/.style={unbreakable}]
\begin{conceptbox}[title={\faWater\ HDBSCAN}, unbreakable]
\small Un point isolé au milieu de nulle part, HDBSCAN le \textbf{laisse de côté} ($-1$), parce qu'il
n'est pas dans une zone dense. Il cherche \textbf{où c'est serré} et accepte de ne pas tout ranger.
\end{conceptbox}
\begin{conceptbox}[title={\faObjectGroup\ Ward}, unbreakable]
\small Le même point isolé, Ward le \textbf{colle de force} au groupe le moins loin, parce qu'il doit
tout ranger. Il \textbf{range tout} et cherche des paquets ronds et équilibrés.
\end{conceptbox}
\end{tcbraster}

\begin{lire}[unbreakable]
\textbf{Et pour les incidents ?} Un incident est une \textbf{zone dense} au milieu de plaintes variées :
c'est pour ça que HDBSCAN est le premier choix. Ward n'est que le plan B, quand les plaintes ne forment pas
d'îles nettes mais un \textbf{continuum}. Là, une méthode de densité ne trouve rien, ou un seul bloc, alors
que Ward découpe quand même quelque chose d'exploitable, avec les mêmes règles d'admission derrière.\par\smallskip
Donc Ward n'est pas \og ajouté après HDBSCAN \fg{} : il est essayé après lui, \textbf{à sa place},
uniquement si HDBSCAN a échoué sur toute la grille. Dans l'exécution normale, \textbf{Ward ne tourne pas
du tout}.
\end{lire}

% Annexe : la stabilité et l'ARI expliqués simplement (exercice 3.5).

\section{La stabilité et l'ARI, comme dans une classe (exercice 3.5)}

\subsection{Le test : refaire les équipes avec des absents}

Une maîtresse range ses élèves en équipes selon leurs goûts. Un jour, quelques élèves sont absents : elle
refait les équipes avec ceux qui restent. Si les amis qui étaient ensemble se retrouvent encore ensemble,
et ceux qui étaient séparés restent séparés, sa façon de faire les équipes est \textbf{fiable}.

Le notebook fait pareil avec les plaintes :
\begin{center}
\begin{tikzpicture}[
  b/.style={draw=#1, fill=#1!10, rounded corners=5pt, align=center, font=\small, text width=3.9cm, minimum height=1.6cm},
  fl/.style={-{Stealth[length=3mm]}, very thick, color=slate}]
\node[b=slate] (a) {{\Large\faDatabase}\\toutes les plaintes\\\textbf{rangement complet}};
\node[b=accentblue, right=1.1cm of a] (b) {{\Large\faRandom}\\on tire \textbf{90 \%} des plaintes\\au hasard};
\node[b=sectioncolor, right=1.1cm of b] (c) {{\Large\faRedo}\\on refait \textbf{tout} le rangement\\sur ces 90 \%};
\draw[fl] (a) -- (b); \draw[fl] (b) -- (c);
\node[b=darkblue, below=0.7cm of b, text width=8.6cm, minimum height=1.1cm] (d) {{\faBalanceScale}\ \ on compare les deux rangements \textbf{sur les mêmes plaintes} : c'est l'ARI};
\draw[fl] (c.south) |- (d.east);
\draw[fl] (a.south) |- (d.west);
\end{tikzpicture}
\end{center}

\begin{tcbraster}[raster columns=2, raster equal height, raster column skip=8pt, raster every box/.style={unbreakable}]
\begin{conceptbox}[title={\faRedo\ Trois fois}, unbreakable]
\small On recommence trois fois. Chaque tirage repart de la base \textbf{complète} : on reste toujours à
90~\%, on ne descend pas à 80 puis 70~\%. Seules les plaintes tirées changent d'une fois à l'autre.
\end{conceptbox}
\begin{conceptbox}[title={\faDiceThree\ Un hasard fixé}, unbreakable]
\small Les tirages utilisent les graines 1, 2 et 3 : si on relance le notebook, on retombe sur les trois
mêmes tirages, donc sur les mêmes résultats.
\end{conceptbox}
\end{tcbraster}

\subsection{Comment on compte : des paires d'élèves}

On ne compare pas les \textbf{numéros} des équipes : l'équipe \og 1 \fg{} d'un jour peut s'appeler
\og 2 \fg{} le lendemain, ça ne veut rien dire. On regarde des \textbf{paires} : pour chaque paire d'élèves,
les deux rangements sont-ils d'accord pour dire \og ensemble \fg{} ou \og séparés \fg{} ?

\begin{center}
\begin{tikzpicture}[el/.style={circle, draw=#1, fill=#1!15, minimum size=0.8cm, font=\small\bfseries}]
\node[font=\small\bfseries, color=darkblue] at (1.6,2.3) {lundi};
\draw[accentblue, thick, rounded corners=8pt] (-0.1,0.6) rectangle (1.5,1.6);
\draw[accentblue, thick, rounded corners=8pt] (1.7,0.6) rectangle (3.3,1.6);
\node[el=accentblue] at (0.3,1.1) {A}; \node[el=accentblue] at (1.1,1.1) {B};
\node[el=ambre] at (2.1,1.1) {C}; \node[el=ambre] at (2.9,1.1) {D};
\begin{scope}[xshift=5cm]
\node[font=\small\bfseries, color=darkblue] at (1.6,2.3) {mardi};
\draw[accentblue, thick, rounded corners=8pt] (-0.1,0.6) rectangle (1.5,1.6);
\draw[accentblue, thick, rounded corners=8pt] (1.7,0.6) rectangle (2.5,1.6);
\draw[accentblue, thick, rounded corners=8pt] (2.6,0.6) rectangle (3.4,1.6);
\node[el=accentblue] at (0.3,1.1) {A}; \node[el=accentblue] at (1.1,1.1) {B};
\node[el=ambre] at (2.1,1.1) {C}; \node[el=ambre] at (3.0,1.1) {D};
\end{scope}
\node[anchor=north west, font=\small, text width=6.2cm] at (9.6,2.4) {
\begin{tabular}{@{}lll@{}}
\textbf{paire} & \textbf{lundi, mardi} & \\
A et B & ensemble, ensemble & \textcolor{accentblue}{\faCheck} \\
A et C & séparés, séparés & \textcolor{accentblue}{\faCheck} \\
A et D & séparés, séparés & \textcolor{accentblue}{\faCheck} \\
B et C & séparés, séparés & \textcolor{accentblue}{\faCheck} \\
B et D & séparés, séparés & \textcolor{accentblue}{\faCheck} \\
C et D & ensemble, séparés & \textcolor{sectioncolor}{\faTimes} \\
\end{tabular}};
\end{tikzpicture}
\end{center}
{\small\color{slate}\textit{Exemple inventé : 4 élèves, 6 paires ; les deux jours sont d'accord sur 5 paires sur 6.}}

\subsection{Pourquoi \og ajusté \fg{} ?}

Compter les paires d'accord donne l'\textbf{indice de Rand}. Mais même deux rangements faits \textbf{au
hasard} tombent souvent d'accord, surtout sur \og séparés \fg, parce que la plupart des paires sont séparées
dans n'importe quel rangement. L'indice de Rand \textbf{ajusté} (ARI) retire cette part due au hasard :

\begin{center}
\begin{tikzpicture}[x=1.35cm]
\shade[left color=palblue, right color=accentblue!80] (0,0) rectangle (10,0.45);
\draw[slate] (0,0) rectangle (10,0.45);
\foreach \x/\v in {0/0, 5/0{,}5, 10/1} { \draw[slate] (\x,0) -- (\x,-0.15); \node[below, font=\small] at (\x,-0.15) {\v}; }
\node[above, font=\footnotesize, text width=3.4cm, align=center] at (0.9,0.5) {pas mieux que le hasard};
\node[above, font=\footnotesize\bfseries, color=sectioncolor, text width=3.4cm, align=center] at (5,0.5) {seuil du notebook};
\node[above, font=\footnotesize, text width=3.4cm, align=center] at (9.1,0.5) {rangements identiques};
\draw[very thick, dashed, color=sectioncolor] (5,-0.05) -- (5,0.5);
\end{tikzpicture}
\end{center}

\doisvoir{trois valeurs d'ARI, une par tirage, toutes au-dessus de 0,5.}

\begin{piege}
\textbf{Stable ne veut pas dire juste.} Mettre tout le monde dans une seule équipe donne un ARI parfait
à chaque fois, et ne range rien. C'est pourquoi l'exercice 3.3 écarte d'abord les rangements ratés, et ne
regarde la stabilité qu'ensuite.
\end{piege}

